% chapters/05_fazit.tex
% Fuenftes Kapitel: Fazit und Ausblick
% Beantwortet die Forschungsfragen, reflektiert die Arbeit und gibt Ausblick.

\chapter{Fazit und Ausblick}
\label{chap:fazit}

\section{Zusammenfassung der Ergebnisse}
\label{sec:zusammenfassung}

Ziel dieser Arbeit war es, [Zielsetzung aus Kapitel~\ref{sec:zielsetzung}
wiederholen]. Ausgehend von den in Kapitel~\ref{chap:ergebnisse} dargestellten
Befunden lassen sich die Forschungsfragen wie folgt beantworten:

\textbf{Forschungsfrage 1:} \ldots

\textbf{Forschungsfrage 2:} \ldots

\textbf{Forschungsfrage 3:} \ldots

\section{Kritische Reflexion}
\label{sec:reflexion}

Die Arbeit liefert [Beitrag benennen]. Dabei sind jedoch die in
Abschnitt~\ref{sec:limitationen} genannten Einschraenkungen zu beachten.
Rueckblickend haette [alternatives Vorgehen] zu schaerferen Ergebnissen
fuehren koennen.

\section{Ausblick}
\label{sec:ausblick}

Aus den Ergebnissen dieser Arbeit ergeben sich mehrere Ansatzpunkte fuer
weiterführende Forschung:

\begin{itemize}
  \item \textbf{Erweiterung der Stichprobe:} Eine breitere Erhebung koennte
        die Generalisierbarkeit der Befunde erhoehen.
  \item \textbf{Laengsschnittstudie:} Ein laengerer Beobachtungszeitraum
        erlaubt Aussagen ueber zeitliche Entwicklungen.
  \item \textbf{Weiterfuehrende Fragestellung:} [Konkrete Forschungsfrage,
        die sich aus dieser Arbeit ergibt] bietet Potenzial fuer
        Anschlussarbeiten.
\end{itemize}
